%% SINH TU figures/make_figures.py -- KHONG sua tay.
\newcommand{\capSizeRatio}{Lợi thế kích thước của EDHOC so với DTLS 1.3 trên chín bộ tham số NIST. Đường đứt: tỉ số cổ điển ở chế độ chữ ký (4.14$\times$). Đường chấm: chế độ Static DH (8.85$\times$), chế độ này KHÔNG có bản hậu lượng tử vì method 1--3 của RFC 9528 dựa trên Diffie--Hellman tĩnh.}
\newcommand{\capExchanges}{Số lượt trao đổi theo kích thước bản tin bắt tay. EDHOC trên CoAP dùng block-wise lock-step nên số lượt nở tuyến tính; DTLS 1.3 truyền theo flight nên giữ nguyên 2 lượt bất kể kích thước [RFC 9147]. Vòng tròn rỗng là SỐ ĐO trên một cài đặt CoAP độc lập, khớp mô hình 7/7. Con số của DTLS là trích chuẩn, không đo ở đây.}
\newcommand{\capBlocksize}{Quét toàn dải cỡ khối mà RFC 7959 cho phép, với bắt tay ML-KEM-768 + ML-DSA-65. Sao là điểm tối ưu của từng đường. Vùng tô và vạch dọc ở 64 B là TRẦN BIÊN DỊCH của RIOT (\texttt{CONFIG\_NANOCOAP\_BLOCK\_SIZE\_MAX}); cỡ khối bên phải vạch không dùng được nếu không sửa firmware. Hai bảng có điểm tối ưu KHÁC NHAU, nên không tồn tại một cấu hình tốt nhất chung cho cả độ trễ lẫn năng lượng.}
\newcommand{\capImplementations}{Ba cài đặt DTLS trên đường truyền ràng buộc; vạch chấm là payload một khung IEEE 802.15.4 (102 B). Số ô bắt tay thành công: gnutls 17/21; mbedtls 14/14; openssl 9/21. Ba cài đặt bị chặn bởi ba thứ khác nhau, nên kết quả này nói về TỪNG CÀI ĐẶT chứ không về giao thức DTLS.}
\newcommand{\capPerMessage}{Kích thước từng bản tin EDHOC. Đường đứt là payload một khung 802.15.4 (102 B). Cột giữa là SÀN của lập luận: biến thể hậu lượng tử nhẹ nhất có thể hình dung, chỉ ML-KEM-512 với xác thực cổ điển và chứng thư theo tham chiếu. Ngay ở sàn đó, riêng khoá đóng gói đã là 800 B, tức 7.8 lần payload khung, nên block-wise bị kích hoạt ngay ở message\_1 với mọi thiết kế.}
